\documentclass[11pt]{article}
\usepackage[T1]{fontenc}
\usepackage{amsmath,amssymb,booktabs}
\usepackage[margin=1in]{geometry}
\usepackage[hidelinks]{hyperref}
\usepackage{tabularx}
\title{Proof Dossier: Rank-Two Trace Centering}
\author{Project Koios Formal}
\date{}
\begin{document}
\maketitle
\section{Identity and status}
\begin{tabularx}{\textwidth}{@{}l>{\raggedright\arraybackslash}X@{}}
\toprule
Statement ID & \texttt{BG-BLOCH-THM-003} \\
Status & \texttt{proved} \\
Lean declaration & \path{ProjectKoios.BandGauge.Bloch.trace_centeredBlochMatrix_eq_zero} \\
Module & \path{ProjectKoios/BandGauge/Bloch/PeriodicRankTwo.lean} \\
\bottomrule
\end{tabularx}

\section{Exact statement and proof}
For every retained $2\times2$ Bloch Hamiltonian matrix $H(k)$,
\[
 D(k)=H(k)-\frac{\operatorname{tr}H(k)}2 I_2
 \quad\Longrightarrow\quad \operatorname{tr}D(k)=0.
\]
The proof expands the two diagonal entries of the trace and applies commutative-ring
algebra.

\section{Source boundary}
This standard supporting theorem has a \texttt{background-only} relation to the trace decomposition at lines 491--496 of the
Ragasa2026A v3 manuscript. Owner: \path{ksdft2effmass}; immutable commit
\begin{center}\small\texttt{99a46cb59d0c6f27ada7d709f194996ef9619d9f}.\end{center}
It is elementary finite matrix algebra and is not a paper-contribution or novelty claim. It does not construct $H(k)$ from a periodic
potential or establish a Bloch--Floquet theorem.

\section{Trusted base}
Lean and pinned Mathlib report \texttt{propext}, \texttt{Classical.choice}, and
\texttt{Quot.sound}. No admission or scientific axiom is used. The proof and prose were
AI-assisted and kernel-checked. No independent expert review is recorded. Status:
\texttt{proved}; working draft.
\end{document}
